\documentclass[11pt]{article}

\usepackage{amsmath,amssymb,amsthm}
\usepackage[margin=1in]{geometry}
\usepackage[hidelinks]{hyperref}

\newtheorem{theorem}{Theorem}
\newtheorem{lemma}{Lemma}
\newtheorem{proposition}{Proposition}
\newtheorem{definition}{Definition}
\newtheorem{remark}{Remark}

\newcommand{\Pow}{\operatorname{powerTailSum}}
\newcommand{\Primes}{\mathcal P}

\title{Certified Tail Bounds for Prime-Power Threshold Cutoff Sequences}
\author{}
\date{}

\begin{document}
\maketitle

\begin{abstract}
For an integer level \(k\ge 1\), put
\[
  \Pow(x,k)=x+x^2+\cdots+x^{k+1}.
\]
If \(t\ge83\) is prime and \(u\) is the next prime after \(t\), define
the prime cutoff \(C_k(t)\) to be the largest prime \(q\), or \(0\) if none
exists, such that
\[
  \Pow(q,k)<\Pow(u,4).
\]
Thus \(C_4(t)=t\).  This paper isolates the threshold-\(83\)
tail-budget theorem
\[
  \sum_{k\ge m} C_k(t)
  \le
  t-\min(3m+1,2m+21,83)
  \qquad (m\ge13).
\]
The proof has an elementary analytic handoff and a finite exact-prefix
certificate.  The analytic side is written out in full.  The finite side is
packaged as a frozen certificate and an independent deterministic verifier.
The verifier recomputes all finite-prefix rows by exact integer arithmetic;
its acceptance is the computational proof of the finite part.
\end{abstract}

\section{Introduction}

This paper studies a concrete family of prime cutoff sequences.  Given a
prime \(t\), let \(u\) be the next prime.  The level-\(k\) cutoff \(C_k(t)\)
is the largest prime whose level-\(k\) prime-power tail lies below the
level-\(4\) threshold determined by \(u\).  The main point is that, once
\(t\ge83\), the late cutoffs have an explicit threshold-row budget:
\[
  \sum_{k\ge m} C_k(t)
  \le
  t-\min(3m+1,2m+21,83)
  \qquad (m\ge13).
\]

The statement is elementary in form.  It involves only primes, finite
geometric sums, and an explicit finite threshold certificate.  Its proof does not use
the Riemann hypothesis, Robin's criterion, or Lagarias's criterion.  Those
criteria explain why this cutoff sequence arose, but they are not needed for
the theorem as stated.  For that reason the paper first develops the cutoff
sequence on its own terms and only later records the divisor-sum motivation.

The proof splits into two parts.  The first part gives a uniform analytic
bound for all sufficiently large threshold rows.  The second part verifies
the remaining finite prefix exactly.  The large-row handoff is compact enough
to write as integer inequalities; the finite prefix is much larger, so it is
handled by an exhaustive verifier using a frozen threshold certificate.  The
accompanying package records the nineteen handoff rows in
\texttt{certificate.csv}.  The verifier recomputes the finite prefix from those
rows without using exploratory search code.

The contribution is therefore deliberately narrow.  First, the paper isolates
the cutoff sequence as an independent elementary object.  Second, it proves a
uniform analytic handoff reducing all large rows to nineteen integer
inequalities.  Third, it supplies a finite-prefix verifier for the
remaining rows.  This is the shape expected of a computational number theory
result: the mathematics reduces the infinite statement to a finite exhaustive
check controlled by a small threshold certificate, and the certificate can be
checked without trusting the search that found it.

\section{Definitions and main theorem}

\begin{definition}[Power-tail sum]
For an integer \(k\ge1\) and a real number \(x\), define
\[
  \Pow(x,k)=\sum_{j=1}^{k+1}x^j.
\]
When \(x=q\) is prime, \(\Pow(q,k)\) is the integer prime-power tail used
to define the cutoff at level \(k\).
\end{definition}

\begin{definition}[Threshold-\(83\) cutoff sequence]
Let \(t\ge83\) be prime, and let \(u\) be the least prime greater than \(t\).
For \(k\ge1\), define
\[
  C_k(t)=
  \max\bigl(\{0\}\cup
  \{q\in\Primes:\Pow(q,k)<\Pow(u,4)\}\bigr).
\]
We usually write \(C_k\) when \(t\) is fixed.  Since \(\Pow(x,4)\) is
strictly increasing and there is no prime strictly between \(t\) and \(u\),
we have \(C_4(t)=t\).
\end{definition}

\begin{theorem}[Threshold-\(83\) exact tail-budget staircase]
\label{thm:threshold83}
Let \(t\ge83\) be prime, let \(u\) be the next prime after \(t\), and define
\(C_k=C_k(t)\) as above.  Then for every integer \(m\ge13\),
\[
  \sum_{k\ge m} C_k
  \le
  t-\min(3m+1,2m+21,83).
\]
\end{theorem}

\begin{remark}
The theorem is certificate-backed by the accompanying package.  The proof
below prints the analytic handoff and specifies the threshold certificate
precisely.  The finite prefix is not printed row by row and is not claimed to
be a single self-contained Lean theorem; instead,
\nolinkurl{verify_tail_certificate.cpp} recomputes the finite prefix
exhaustively from the frozen certificate by exact integer arithmetic.
\end{remark}

\section{Proof strategy}

The proof begins with a one-line consequence of the cutoff definition.  If
\(C_k>0\), then
\[
  \Pow(C_k,k)<\Pow(u,4).
\]
The lower bound \(\Pow(q,k)>q^{k+1}\), the upper bound
\(\Pow(u,4)\le2u^5\), and Bertrand's postulate \(u<2t\) imply
\[
  C_k^{k+1}<64t^5.
\]
Therefore
\[
  C_k<(64t^5)^{1/(k+1)}.
\]
There are also only \(O(\log t)\) nonzero levels, because \(C_k>0\) forces
\(2^{k+1}<64t^5\).  Summing the uniform envelope over \(k\ge m\) gives
\[
  \sum_{k\ge m}C_k
  \le
  64^{1/(m+1)}(6+5\log_2 t)t^{5/(m+1)}.
\]
For \(t\ge256\), the elementary logarithmic envelope
\[
  6+5\log_2 t\le 36t^{1/4}
\]
turns this into a pure power inequality.  For each \(13\le m\le31\), it
then suffices to verify one integer handoff row:
\[
  36^{4a}64^4T^{20+a}\le (T-g_m)^{4a},
  \qquad a=m+1,\quad g_m=\min(3m+1,2m+21,83).
\]
The exponent in the corresponding ratio is less than \(1\), so once the
inequality holds at the handoff \(T\), it holds for every larger \(t\).  The
finite range \(83\le t<T\) is checked exactly from the cutoff definition.
For \(m\ge31\), monotonicity of the tail in \(m\) reduces the result to the
case \(m=31\), where the staircase has already reached the constant gap
\(83\).

\section{Proof of the theorem}

\begin{lemma}[Active cutoff envelope]
\label{lem:active-envelope}
Let \(t\ge83\) be prime and \(u\) the next prime.  If \(C_k(t)>0\), then
\[
  C_k(t)^{k+1}<64t^5.
\]
\end{lemma}

\begin{proof}
By definition,
\[
  \Pow(C_k,k)<\Pow(u,4).
\]
For \(q\ge2\), \(\Pow(q,k)>q^{k+1}\).  Also
\[
  \Pow(u,4)=u+u^2+u^3+u^4+u^5\le2u^5
\]
for \(u\ge2\).  Bertrand's postulate gives \(u<2t\).  Hence
\[
  C_k^{k+1}<2u^5<2(2t)^5=64t^5.
\]
\end{proof}

\begin{lemma}[Smooth tail majorant]
\label{lem:smooth-tail}
For every \(m\ge1\),
\[
  \sum_{k\ge m}C_k(t)
  \le
  64^{1/(m+1)}(6+5\log_2 t)t^{5/(m+1)}.
\]
\end{lemma}

\begin{proof}
If \(C_k>0\), then \(C_k\ge2\), so Lemma~\ref{lem:active-envelope} gives
\[
  2^{k+1}<64t^5.
\]
Thus
\[
  k+1<6+5\log_2 t.
\]
For every nonzero term with \(k\ge m\), Lemma~\ref{lem:active-envelope} also
gives
\[
  C_k<(64t^5)^{1/(k+1)}
  \le
  (64t^5)^{1/(m+1)}.
\]
There are fewer than \(6+5\log_2 t\) possible nonzero terms.  Multiplying the
number of possible terms by the largest displayed bound gives the result.
\end{proof}

\begin{lemma}[Quarter-root analytic handoff]
\label{lem:handoff}
Let \(13\le m\le31\), put \(a=m+1\), and put
\[
  g_m=\min(3m+1,2m+21,83).
\]
Assume \(T\ge256\) and
\[
  36^{4a}64^4T^{20+a}\le (T-g_m)^{4a}.
  \tag{1}
\]
Then for every \(t\ge T\),
\[
  \sum_{k\ge m}C_k(t)\le t-g_m.
\]
\end{lemma}

\begin{proof}
For \(t\ge256\), the elementary estimate
\[
  6+5\log_2 t\le36t^{1/4}
\]
holds.  Indeed, at \(t=256\) the right side is \(144\), while the left side
is \(46\), and the difference is increasing for \(t\ge256\) because
\[
  \frac{d}{dt}\left(36t^{1/4}-6-5\log_2 t\right)
  =
  \frac{1}{t}\left(9t^{1/4}-\frac{5}{\log 2}\right)>0.
\]
Combining this estimate with Lemma~\ref{lem:smooth-tail} gives
\[
  \sum_{k\ge m}C_k(t)
  \le
  36\,64^{1/a}t^{5/a+1/4}.
\]
The right side is at most \(t-g_m\) exactly when
\[
  36^{4a}64^4t^{20+a}\le (t-g_m)^{4a}.
  \tag{2}
\]
The ratio of the left side of (2) to the right side is a constant times
\[
  \frac{t^{20+a}}{(t-g_m)^{4a}}.
\]
Equivalently, after taking \(4a\)-th roots, it is enough to show that
\[
  36\,64^{1/a}t^{(20+a)/(4a)}\le t-g_m.
\]
Put \(r=(20+a)/(4a)\).  For \(m\ge13\) we have \(a\ge14\), hence \(0<r<1\).
The function
\[
  \frac{t^r}{t-g_m}
\]
is strictly decreasing for \(t>g_m\), since
\[
  \frac{d}{dt}\log \frac{t^r}{t-g_m}
  =
  \frac{r}{t}-\frac{1}{t-g_m}
  =
  \frac{(r-1)t-rg_m}{t(t-g_m)}<0.
\]
Thus once (1) holds at \(T\), (2) holds for all \(t\ge T\).  The repository
formalizes this integer monotonicity statement as the quarter-root handoff.
The claimed tail bound follows.
\end{proof}

\begin{proposition}[Finite-prefix certificate]
\label{prop:finite-certificate}
For \(13\le m\le31\), define \(g_m=\min(3m+1,2m+21,83)\).  The following
handoff thresholds satisfy the hypotheses of Lemma~\ref{lem:handoff}:
\[
\begin{array}{c|cccccccccc}
m&13&14&15&16&17&18&19&20&21&22\\
\hline
T_m&
19596&10674&6641&4542&3332&2578&2079&1733&1481&1293
\end{array}
\]
\[
\begin{array}{c|ccccccccc}
m&23&24&25&26&27&28&29&30&31\\
\hline
T_m&
1149&1037&947&873&813&763&721&685&654
\end{array}
\]
Moreover, for every prime \(t\) with \(83\le t<T_m\), the exact cutoff
sequence satisfies
\[
  \sum_{k\ge m}C_k(t)\le t-g_m.
\]
\end{proposition}

\begin{proof}
The first assertion is a finite list of integer inequalities of the form
\[
  36^{4a}64^4T_m^{20+a}\le (T_m-g_m)^{4a}.
\]
These inequalities are small enough to be checked directly by integer
arithmetic and are checked by the packaged verifier.  The same handoff rows
are also mirrored by the repository's Lean handoff table.

The second assertion is the finite exact-prefix certificate.  The verifier
recomputes the successor prime \(u\), recomputes each \(C_k(t)\) from the
inequality
\[
  \Pow(C_k,k)<\Pow(u,4),
\]
sums the tail, and checks the displayed inequality.  No floating-point
arithmetic is needed for this finite prefix.  In the present paper this
finite row list is not printed; Section~\ref{sec:certificate} specifies the
packaged certificate and verifier that audit it independently.
\end{proof}

\begin{proof}[Proof of Theorem~\ref{thm:threshold83}]
First suppose \(13\le m\le31\), and put
\[
  g_m=\min(3m+1,2m+21,83).
\]
If \(83\le t<T_m\), Proposition~\ref{prop:finite-certificate} gives
\[
  \sum_{k\ge m}C_k(t)\le t-g_m.
\]
If \(t\ge T_m\), Lemma~\ref{lem:handoff} gives the same inequality from the
integer handoff row.  Hence the theorem holds for \(13\le m\le31\).

If \(m\ge31\), then
\[
  \sum_{k\ge m}C_k(t)\le \sum_{k\ge31}C_k(t)\le t-83.
\]
Since
\[
  \min(3m+1,2m+21,83)=83
  \qquad(m\ge31),
\]
this is exactly the asserted bound.
\end{proof}

\section{Certificate and verifier architecture}
\label{sec:certificate}

The theorem is well suited to referee-grade computational verification because
the threshold certificate can be made small in concept and deterministic in
execution.  The verifier should not depend on exploratory search code.  It
should accept only a frozen certificate and either accept or reject it by
direct arithmetic.

The computation package consists of four audit files:
\begin{center}
\begin{tabular}{ll}
certificate & \nolinkurl{certificate.csv}\\
verifier & \nolinkurl{verify_tail_certificate.cpp}\\
instructions & \nolinkurl{BUILD_AND_VERIFY.txt}\\
hashes & \nolinkurl{SHA256SUMS}
\end{tabular}
\end{center}
The finite part is used in the following acceptance form: the verifier
\nolinkurl{verify_tail_certificate}, compiled with the stated compiler version
and invoked by the stated command line, accepts the threshold-certificate file
with the stated SHA256 digest; acceptance implies
Proposition~\ref{prop:finite-certificate}.  This is stronger than citing an
exploratory repository script, because the proof obligation becomes the
verifier's acceptance theorem.

The version-\(1.0.0\) verification artifact is archived at Zenodo.  Its
versioned DOI is recorded in the repository metadata.  The artifact uses split
licensing; its file-level license map is in \texttt{LICENSE.md}.

The threshold certificate should contain the following data.

\begin{enumerate}
\item The nineteen handoff rows \((m,T_m,g_m)\) for \(13\le m\le31\).
\item A file hash for the certificate, a hash for the verifier source, and
the exact command line used to reproduce the acceptance result.
\end{enumerate}

The verifier should perform the following checks.

\begin{enumerate}
\item Recompute primality of every \(t\) and recompute that \(u\) is the least
prime greater than \(t\).
\item Recompute \(\Pow(u,4)\) and each cutoff \(C_k(t)\) using integer
arithmetic.  For a positive cutoff \(C_k\), check
\[
  \Pow(C_k,k)<\Pow(u,4),
\]
and either \(C_k=0\) or the next prime after \(C_k\) fails the same inequality.
\item Recompute the tail sum for every prime \(t\) with \(83\le t<T_m\) and
check
\[
  \sum_{k\ge m}C_k(t)\le t-g_m.
\]
\item Check the handoff inequalities
\[
  36^{4a}64^4T_m^{20+a}\le (T_m-g_m)^{4a},
  \qquad a=m+1.
\]
\item Check the formula \(g_m=\min(3m+1,2m+21,83)\) for every listed \(m\).
\end{enumerate}

The packaged verifier is deliberately independent of the exploratory scripts
that found the thresholds.  It uses its own prime sieve and unsigned
base-\(10^9\) big-integer arithmetic, and it uses no floating-point
arithmetic.

The archived artifact keeps the finite audit path deliberately small: the
standalone C++ verifier is the shipped executable specification of the finite
prefix.  No exploratory threshold-generation script, floating-point estimate,
or external proof-assistant project is needed to check the certificate in the
archive.

The minimal verifier is only a few logical steps.  In pseudocode, for each
listed pair \((m,t)\) it should do the following:
\[
\begin{array}{l}
\text{assert } t \text{ is prime and } 83\le t<T_m,\\
u\leftarrow\text{least prime greater than }t,\\
P_4\leftarrow \Pow(u,4),\\
\text{for } k=m,m+1,\ldots\text{ while }\Pow(2,k)<P_4:\\
\qquad C_k\leftarrow\max\{0\}\cup\{q\in\Primes:\Pow(q,k)<P_4\},\\
\text{assert } \sum_{k\ge m}C_k\le t-g_m.
\end{array}
\]
The loop is finite because \(C_k>0\) implies \(\Pow(2,k)<P_4\).  This
pseudocode is intentionally independent of any heuristic row generation.  It
is the verification contract a referee would need to trust.

\section{Motivation from divisor-sum extremal analysis}

The cutoff sequence above arose from divisor-sum extremal computations, but
the theorem itself is independent of that origin.  In Robin- and
Lagarias-type analyses one often studies integers
\[
  n=\prod_{p\le P}p^{a_p}
\]
through the normalized divisor factors
\[
  A_p(a)=1+\frac1p+\cdots+\frac1{p^a}.
\]
A one-step marginal comparison has the form
\[
  \frac{A_p(a+1)}{A_p(a)}
  =
  1+\frac{p-1}{p(p^{a+1}-1)}.
\]
After clearing denominators, this naturally produces sums of prime powers.
For example, the integer quantity
\[
  p+p^2+\cdots+p^{a+1}
\]
is the reciprocal scale of the marginal gain.  Cutoffs defined by inequalities
between such power tails are therefore a natural bookkeeping device for
organizing which prime powers remain active in an extremal envelope.

Classically, Robin's criterion and Lagarias's criterion connect divisor-sum
inequalities with the Riemann hypothesis, and the superabundant and
colossally abundant frameworks of Alaoglu and Erdos explain why exponent
windows and marginal comparisons appear.  The present theorem should not be
read as an RH result.  Its role is narrower: it gives an explicit certified
tail budget for one family of prime-power threshold cutoffs that arose inside
that broader extremal-number analysis.

\section{Significance and limitations}

The useful feature of Theorem~\ref{thm:threshold83} is its exactness.  The
right side is not merely asymptotic; it is a concrete staircase with sharp
threshold-row behavior in the available data.  The theorem also isolates a
clean computational object: a finite prefix plus a transparent analytic
handoff.

The main limitation is formalization status.  The analytic handoff is
elementary and the finite prefix is certified by the packaged verifier, but the
submitted artifact is still a computational proof package rather than one small
end-to-end proof-assistant theorem named exactly as
Theorem~\ref{thm:threshold83}.

A second limitation is scope.  The theorem concerns one threshold, namely
\(83\), and one power-tail cutoff model.  It does not by itself prove a Robin,
Lagarias, or Riemann-hypothesis statement.  Its value is as a reusable explicit
tail bound inside extremal divisor-sum analysis and as a standalone example of
a certified prime-power cutoff staircase.

\begin{thebibliography}{9}

\bibitem{Lagarias2002}
J. C. Lagarias,
\emph{An Elementary Problem Equivalent to the Riemann Hypothesis},
American Mathematical Monthly 109 (2002), 534--543;
arXiv:math/0008177.

\bibitem{Robin1984}
G. Robin,
\emph{Grandes valeurs de la fonction somme des diviseurs et hypothese de
Riemann},
Journal de Mathematiques Pures et Appliquees 63 (1984), 187--213.

\bibitem{AlaogluErdos1944}
L. Alaoglu and P. Erdos,
\emph{On highly composite and similar numbers},
Transactions of the American Mathematical Society 56 (1944), 448--469.

\bibitem{GrahamKnuthPatashnik1994}
R. L. Graham, D. E. Knuth, and O. Patashnik,
\emph{Concrete Mathematics: A Foundation for Computer Science},
2nd edition, Addison-Wesley, 1994.

\end{thebibliography}

\end{document}
